% Chapter fragment for repository integration.
% Assumes theorem environments and ClaimStatus macros from main.tex.
% No new macros are introduced here.

\chapter{The affine half-space BV theorem and its transport to Virasoro and \texorpdfstring{$W_3$}{W3}}
\label{chap:affine-half-space-bv}

\begin{remark}[Position in the theorem graph]
This chapter proves the half-space deformation-quantization programme for
Poisson vertex algebras.  The affine case
(Theorem~\ref{thm:affine-half-space-bv}) is the explicit one-loop-exact
special case, displaying the essential passage
\[
\text{bulk BV theory} \longrightarrow \text{boundary VOA} \longrightarrow \text{Drinfeld--Sokolov reduction}
\]
in full detail.  The general case
(Theorem~\ref{thm:general-half-space-bv}) is proved via the doubling
construction (\S\ref{thm:doubling-rwi}), which embeds the half-space
problem into the standard whole-space Kontsevich machinery.  The
central point of the affine case is that the affine Poisson--vertex
sigma model on $\mathbb C \times \mathbb R_{\ge 0}$ \emph{is} the mixed
BF/Chern--Simons theory, and the half-space BV quantization theorem of
Gwilliam--Williams solves it to all loop orders with a single formal
level parameter.
\end{remark}

\begin{summary}[What is proved here]
The main theorem of this chapter is the following.

\medskip

\emph{For a semisimple Lie algebra $\mathfrak g$, the affine Poisson--vertex sigma model of
Khan--Zeng on the half-space $X = \mathbb C \times \mathbb R_{\ge 0}$ with Neumann boundary
condition admits a renormalized BV quantization to all orders in $\hbar$; in fact the quantization
is one-loop exact.  The resulting boundary vertex algebra is the affine Kac--Moody vertex algebra
$V^{K_{\mathrm{eff}}}(\mathfrak g)$ at a uniquely shifted formal level
$K_{\mathrm{eff}}(k,\hbar)=k+\hbar\,\delta_{\mathfrak g}$, and there are no higher-loop boundary anomaly
channels.  Principal Drinfeld--Sokolov reduction then transports this single formal modulus to the
Virasoro and $W_3$ families.}
\end{summary}

\section{The question in first principles form}

\begin{problem}[The half-space BV problem]
Let $V$ be a freely generated Poisson vertex algebra, and let
$X = \mathbb C_z \times \mathbb R_{\ge 0,t}$.
The three-dimensional Poisson--vertex sigma model attaches to $V$ a mixed holomorphic-topological
BV theory on $X$ whose Neumann boundary condition keeps the $\Phi$-fields free and sets the
$\eta$-fields to vanish at the boundary.  The sought deformation quantization of $V$ is the boundary
vertex algebra of the quantized half-space theory.

The missing theorem is the existence of a renormalized solution of the quantum master equation on the
half-space, together with control of the boundary anomaly classes at all loop orders.
\end{problem}

\begin{remark}[Why affine is the first theorematic case]
The affine Poisson vertex algebra is the first case where the nonlinear Poisson--vertex sigma model
reduces to a chiral deformation of a mixed BF theory, hence to a class for which a rigorous BV
quantization theorem already exists.  The point of this chapter is that the half-space theorem sought by
Khan--Zeng is already available in this affine case, once the fields and the boundary condition are
identified correctly.
\end{remark}

\section{The affine Poisson--vertex sigma model from first principles}

\begin{setup}[Affine PVA and half-space fields]
Let $\mathfrak g$ be a finite-dimensional semisimple Lie algebra over $\mathbb C$, with structure constants
$f^{ab}{}_c$ and invariant bilinear form $\kappa^{ab}$.
The affine Poisson vertex algebra is
\[
V_{\mathrm{aff}}(\mathfrak g,k) = \mathrm{Sym}\bigl(\mathbb C[\partial] \otimes \mathfrak g\bigr)
\]
with generators $B^a$ and $\lambda$-bracket
\[
\{ B^a{}_{\lambda} B^b \} = f^{ab}{}_c B^c + k\,\kappa^{ab}\lambda.
\]
The corresponding three-dimensional Poisson--vertex sigma model on
$X = \mathbb C \times \mathbb R_{\ge 0}$ has fields
\[
B = B^a t_a \in \Omega^{0,0}(X) \otimes \mathfrak g,
\qquad
A = A_t\,dt + A_{\bar z}\,d\bar z \in \Omega^{0,1}_{\mathrm{mix}}(X) \otimes \mathfrak g,
\]
with classical action
\begin{equation}
\label{eq:affine-kz-action}
S_{\mathrm{aff}}(B,A)
=
\int_X \! dz\,
\left(
\kappa_{ab} B^a (d_t + \bar\partial) A^b
+ \frac{1}{2} f_{abc} B^a A^b A^c
+ \frac{k}{2}\kappa_{ab} A^a \partial A^b
\right).
\end{equation}
The Neumann boundary condition is
\begin{equation}
\label{eq:kz-neumann}
i^*A = 0,
\qquad
B\ \text{free on }\partial X,
\end{equation}
where $i\colon \partial X \hookrightarrow X$ is the boundary inclusion.
\end{setup}

\begin{remark}[Affine as Chern--Simons for $k\neq 0$]
When $k\neq 0$, the field redefinition
\[
\widetilde A = A + \frac{1}{k}\kappa^{-1}(B)\,dz
\]
transforms the bulk action to analytically continued Chern--Simons theory.
This is the classical reason the affine example is fully topological once the inner Sugawara/Virasoro
field is present.
\end{remark}

\section{Identification with mixed BF/Chern--Simons}

\begin{proposition}[Dictionary with mixed BF theory]\label{prop:affine-mixed-bf-dictionary}
\ClaimStatusProvedHere.
Let
\[
\alpha := A \in \Omega^{0,\bullet}(\mathbb C)\widehat\otimes \Omega^{\bullet}(\mathbb R_{\ge 0})[1] \otimes \mathfrak g,
\qquad
\beta := B\,dz \in \Omega^{1,\bullet}(\mathbb C)\widehat\otimes \Omega^{\bullet}(\mathbb R_{\ge 0}) \otimes \mathfrak g.
\]
Then the affine action \eqref{eq:affine-kz-action} is exactly
\begin{equation}
\label{eq:smix-plus-ics}
S_{\mathrm{aff}} = S_{\mathrm{mix}} + k\,I_{\mathrm{CS}},
\end{equation}
where
\begin{equation}
\label{eq:smix-def}
S_{\mathrm{mix}}(\alpha,\beta)
=
\int_X \langle \beta,(\bar\partial+d_t)\alpha\rangle
+ \frac{1}{2}\int_X \langle \beta,[\alpha,\alpha]\rangle,
\end{equation}
and
\begin{equation}
\label{eq:ics-def}
I_{\mathrm{CS}}(\alpha)
=
\frac{1}{2}\int_X \langle \alpha,\partial\alpha\rangle.
\end{equation}
In particular, the affine Poisson--vertex sigma model is a chiral deformation of mixed BF theory.
\end{proposition}

\begin{proof}
The field $B$ is a scalar, so $\beta = B\,dz$ has bidegree $(1,\bullet)$ in the mixed holomorphic-topological
splitting.  The kinetic term in \eqref{eq:affine-kz-action} becomes
\[
\int_X \kappa_{ab} B^a (d_t+\bar\partial)A^b\,dz
=
\int_X \langle \beta,(\bar\partial+d_t)\alpha\rangle.
\]
Likewise the cubic term rewrites as
\[
\frac{1}{2}\int_X f_{abc} B^a A^b A^c\,dz
=
\frac{1}{2}\int_X \langle \beta,[\alpha,\alpha]\rangle.
\]
Finally,
\[
\frac{k}{2}\int_X \kappa_{ab} A^a \partial A^b\,dz
=
\frac{k}{2}\int_X \langle \alpha,\partial\alpha\rangle
=
 k\,I_{\mathrm{CS}}(\alpha).
\]
This is precisely the mixed BF theory together with its chiral Chern--Simons deformation.
\end{proof}

\begin{remark}[Combinatorial consequence]
The chiral nature of the deformation is decisive.
In mixed BF theory, every propagator runs from an $\alpha$-leg to a $\beta$-leg, and every chiral interaction
has at most one $\beta$-leg.  Hence the only connected Feynman graphs that can occur are trees,
single vertices, or one-loop graphs.  This is the combinatorial source of one-loop exactness.
\end{remark}

\section{The boundary condition as a derived Lagrangian boundary condition}

\begin{definition}[Chiral boundary condition in mixed BF language]
Let $\sigma\colon \mathbb R \to \mathbb R$ be reflection, $\sigma(t)=-t$.
Write $A^{\mathrm{odd}} \subset \Omega^{0,\bullet}(\mathbb C)\widehat\otimes\Omega^{\bullet}(\mathbb R)$
for the $(-1)$-eigenspace under pullback by $\sigma$.
Concretely, its $0$-form part consists of functions $f(t)$ with $f(0)=0$, and its $1$-form part consists
of one-forms $g(t)\,dt$ whose pullback to the boundary vanishes because $i^*(dt)=0$.
The chiral boundary condition is the restriction
\[
\alpha \in A^{\mathrm{odd}}\otimes \mathfrak g[1],
\qquad
\beta\ \text{free}.
\]
\end{definition}

\begin{proposition}[Equivalence of boundary conditions]\label{prop:boundary-condition-equivalence}
\ClaimStatusProvedHere.
The Neumann boundary condition \eqref{eq:kz-neumann} for the affine Poisson--vertex sigma model is
implemented, in the mixed BF/Chern--Simons presentation, by the chiral boundary condition above.
More precisely, the odd-reflection model is a strict model for the derived fiber product imposing
$i^*\alpha=0$.
\end{proposition}

\begin{proof}
If $\alpha \in A^{\mathrm{odd}}\otimes \mathfrak g[1]$, then its $0$-form part is an odd function of $t$ and hence
vanishes at $t=0$; its $dt$-part survives as an odd coefficient times $dt$, but its pullback to the boundary
also vanishes because $i^*(dt)=0$.  Therefore $i^*\alpha=0$, which is exactly the Khan--Zeng condition
$i^*A=0$.

Conversely, the BV formalism does not impose boundary conditions by naive degreewise kernels alone; one must
produce a Lagrangian subcomplex modelling the derived fiber product of fields with the zero section of the boundary
restriction map.  The odd-reflection model does exactly this.  Because the restriction map to the boundary is
levelwise surjective, the strict fiber product coincides with the derived one.  Hence the chiral boundary condition
is a BV-correct model of the Neumann boundary condition.
\end{proof}

\begin{remark}[First-principles content]
This proposition is the conceptual hinge of the chapter.
The half-space problem posed for the affine Poisson--vertex sigma model is not a new analytic problem;
it is the already-solved half-space problem for mixed BF/Chern--Simons, expressed in different fields.
\end{remark}

\section{The all-loops half-space BV theorem in the affine case}

\begin{maintheorem}[Affine half-space BV quantization]\label{thm:affine-half-space-bv}
\ClaimStatusProvedHere.
Let $\mathfrak g$ be semisimple and let $k\in\mathbb C$.
Consider the affine Poisson--vertex sigma model on
$X=\mathbb C\times\mathbb R_{\ge 0}$ with Neumann boundary condition.
Then the following hold.
\begin{enumerate}
    \item There exists a renormalized perturbative BV quantization solving the quantum master equation to all orders in $\hbar$.
    \item The quantization is one-loop exact: the renormalized effective action receives no contributions from loop order $\ge 2$.
    \item The only quantum correction is a local one-loop functional proportional to
    \begin{equation}
    \label{eq:one-loop-boundary-shift-functional}
    \hbar\int_X \operatorname{Tr}^{\mathrm{ad}}_{\mathfrak g}(\alpha\,\partial\alpha).
    \end{equation}
    \item Consequently the quantum theory is equivalent to the affine theory at a uniquely shifted formal level
    \begin{equation}
    \label{eq:keff-def}
    K_{\mathrm{eff}}(k,\hbar) = k + \hbar\,\delta_{\mathfrak g},
    \end{equation}
    where $\delta_{\mathfrak g}$ is proportional to the invariant bilinear form on $\mathfrak g$.
\end{enumerate}
\end{maintheorem}

\begin{proof}
By Proposition \ref{prop:affine-mixed-bf-dictionary}, the affine theory is a chiral deformation of mixed BF
by the Chern--Simons functional $kI_{\mathrm{CS}}$.  By Proposition \ref{prop:boundary-condition-equivalence},
its Neumann boundary condition is precisely the chiral boundary condition in the mixed BF presentation.

The half-space BV theorem for mixed BF/Chern--Simons now applies directly: for any chiral deformation of mixed
BF theory on $\mathbb C\times\mathbb R_{\ge 0}$ in the presence of the chiral boundary condition, there is no anomaly
and there exists a one-loop exact quantization.  Specializing to the affine deformation gives items (1) and (2).
The same theorem identifies the one-loop correction with a local functional proportional to
\eqref{eq:one-loop-boundary-shift-functional}, giving item (3).

For semisimple $\mathfrak g$, translation-invariant cotangent quantizations of mixed BF theory form a torsor for
$\hbar\,\mathrm{Sym}^2(\mathfrak g)^{\mathfrak g}$.  Since $\mathfrak g$ is semisimple, this space is one-dimensional on each
simple summand and generated by the invariant bilinear form.  Thus the one-loop correction has no freedom other
than a level shift, yielding \eqref{eq:keff-def}.  Since there are no loops beyond one, the formal level has no
higher $\hbar$-terms.  This proves item (4).
\end{proof}

\begin{corollary}[The affine case of the Khan--Zeng half-space problem]\label{cor:affine-kz-solved}
\ClaimStatusProvedHere.
For the affine Poisson vertex algebra, the half-space deformation quantization problem posed by Khan--Zeng is
unobstructed.  The obstruction theory terminates at one loop, and that one-loop term is absorbed by the unique
formal level shift $k\mapsto K_{\mathrm{eff}}(k,\hbar)$.
\end{corollary}

\begin{proof}
Theorem \ref{thm:affine-half-space-bv} gives a renormalized BV quantization to all orders.  One-loop exactness
kills every possible higher-loop obstruction.  The remaining one-loop functional is exactly the level-shift term,
which lies in the one-dimensional quantization torsor and hence is absorbed into the renormalized level.
\end{proof}

\begin{remark}[Physical normalization]
The theorem leaves the coefficient $\delta_{\mathfrak g}$ abstract.  In standard vertex-algebra and
Chern--Simons/WZW normalizations one expects
$\delta_{\mathfrak g}=-h^{\vee}$ on each simple summand, but the analytic input above only proves
that the one-loop term is proportional to the invariant form, not the exact proportionality constant.
The theorematic content of the half-space quantization does not require pinning down that scalar.
\end{remark}

\section{The boundary obstruction line}

\begin{definition}[Boundary jet obstruction complex]\label{def:boundary-jet-obstruction}
Let
\[
\mathcal A_{\partial,\mathrm{aff}}
=
\mathbb C\bigl[J^a_{(m)},c^a_{(m)} \,;\, a=1,\dots,\dim\mathfrak g,\ m\ge 0\bigr]
\]
be the boundary jet algebra generated by the current jets
$J^a_{(m)}=\partial^m J^a$ and ghost jets $c^a_{(m)}=\partial^m c^a$.
Set
\[
\operatorname{Obs}^{\bullet}_{\partial,\mathrm{aff}}
=
\mathcal A_{\partial,\mathrm{aff}}/\partial \mathcal A_{\partial,\mathrm{aff}},
\]
graded by ghost number, and define the BRST differential by
\[
\delta J^a_{(m)}
=
K\,c^a_{(m+1)}
+
f^a{}_{bc}\sum_{r=0}^{m}\binom{m}{r}J^b_{(r)}c^c_{(m-r)},
\]
\[
\delta c^a_{(m)}
=
-\frac12 f^a{}_{bc}\sum_{r=0}^{m}\binom{m}{r}c^b_{(r)}c^c_{(m-r)}.
\]
Degree $0$ elements are local counterterms, degree $1$ elements are local boundary anomalies, and degree $2$
elements are obstructions to removing those anomalies.
\end{definition}

\begin{remark}[Current-preserving truncation]
The two-term complex $C^{\bullet}_{\mathrm{KM}}$ introduced below is the weight-one, singular-OPE truncation of
Definition \ref{def:boundary-jet-obstruction}.  It is the correct deformation complex for the affine
current sector, i.e.\ for deformations that preserve a weight-one current generator and keep the singular OPE
linear in that current.
\end{remark}

\begin{construction}[The current-preserving boundary deformation complex]\label{cons:boundary-km-complex}
Let $J_x(z)$ denote the boundary current associated to $x\in \mathfrak g$.
A first-order current-preserving singular deformation of the affine OPE is a pair
\[
(\omega,\mu)\in \operatorname{Sym}^2(\mathfrak g^*)\oplus \operatorname{Hom}(\Lambda^2\mathfrak g,\mathfrak g)
\]
with deformed singular OPE
\[
J_x(z)J_y(w)
\sim
\frac{K\kappa(x,y)+\varepsilon\,\omega(x,y)}{(z-w)^2}
+
\frac{J_{[x,y]+\varepsilon\mu(x,y)}(w)}{z-w}.
\]
A first-order field redefinition is $R\in\operatorname{End}(\mathfrak g)$ acting by
$J_x\mapsto J_x+\varepsilon J_{Rx}$.
This yields a two-term complex
\[
C^0_{\mathrm{KM}} = \operatorname{End}(\mathfrak g),
\qquad
C^1_{\mathrm{KM}} = \operatorname{Sym}^2(\mathfrak g^*)\oplus \operatorname{Hom}(\Lambda^2\mathfrak g,\mathfrak g),
\]
with differential
\[
d_0R = (\omega_R,\mu_R),
\]
where
\[
\omega_R(x,y) = K\bigl(\kappa(Rx,y)+\kappa(x,Ry)\bigr),
\qquad
\mu_R(x,y)=R[x,y]-[Rx,y]-[x,Ry].
\]
The linearized Jacobi and invariance equations are
\[
\delta_{\mathrm{CE}}\mu = 0,
\qquad
\omega([x,y],z)-\omega(x,[y,z]) + K\kappa(\mu(x,y),z)-K\kappa(x,\mu(y,z))=0.
\]
\end{construction}

\begin{theorem}[The affine boundary deformation line]\label{thm:affine-boundary-line}
\ClaimStatusProvedHere.
For simple $\mathfrak g$,
\[
H^1(C^{\bullet}_{\mathrm{KM}}) \cong \mathbb C\cdot \Theta,
\qquad
\Theta = [\,(\kappa,0)\,].
\]
Equivalently, every infinitesimal current-preserving singular deformation of the affine boundary OPE is
gauge-equivalent to a unique infinitesimal level shift.
\end{theorem}

\begin{proof}
Let $(\omega,\mu)$ be a cocycle.  The linearized Jacobi identity says that $\mu$ is a Chevalley--Eilenberg
$2$-cocycle with values in the adjoint module.  Since $\mathfrak g$ is simple, Whitehead's lemma gives
$H^2(\mathfrak g;\mathfrak g)=0$, so $\mu=\delta_{\mathrm{CE}}R$ for some $R\in\operatorname{End}(\mathfrak g)$.
Subtracting the coboundary $d_0R$ kills $\mu$, so the cohomology class is represented by $(\omega',0)$.
The second cocycle equation then reduces to invariance of $\omega'$, namely
\[
\omega'([x,y],z)=\omega'(x,[y,z]).
\]
Thus $\omega'$ is an invariant symmetric bilinear form on the simple Lie algebra $\mathfrak g$, hence a scalar
multiple of the Killing form.  Therefore the cohomology is one-dimensional and generated by the class
$[(\kappa,0)]$.
\end{proof}

\begin{lemma}[Pole--weight balance for boundary currents]\label{lem:pole-weight-balance}
\ClaimStatusProvedHere.
In the affine half-space theory, every singular contribution to the boundary current OPE
\[
J^a(z)J^b(w)\sim \sum_{m\ge 1} \frac{{\mathcal O}^{ab}_m(w)}{(z-w)^m}
\]
satisfies
\begin{equation}
\label{eq:pole-weight-balance}
m + \mathrm{wt}({\mathcal O}^{ab}_m)=2,
\end{equation}
where $\mathrm{wt}(J)=1$ and $\mathrm{wt}(\partial)=1$.
\end{lemma}

\begin{proof}
Each external current has weight $1$, so the input weight is $2$.
In the mixed BF/Chern--Simons presentation one assigns
$\mathrm{wt}(\beta)=1$, $\mathrm{wt}(\alpha)=0$, and $\mathrm{wt}(\partial)=1$.
Every bulk interaction then has total holomorphic weight $1$, while the reflected propagator is homogeneous
of degree $-1$ in the boundary separation.  Renormalized local counterterms preserve this scaling because there is
no additional boundary tensor available.  Hence any singular term in the OPE of two weight-one currents has
homogeneous total weight $2$.  A pole of order $m$ contributes weight $m$, so the remaining operator must have
weight $2-m$.
\end{proof}

\begin{theorem}[Current-sector rigidity of the quantum boundary algebra]\label{thm:current-sector-rigidity}
\ClaimStatusProvedHere.
Let $\mathfrak g$ be simple.  The singular OPE of the quantum boundary currents has the exact affine form
\begin{equation}
\label{eq:quantum-affine-boundary-ope}
J^a(z)J^b(w)
\sim
\frac{K_{\mathrm{eff}}(k,\hbar)\,\kappa^{ab}}{(z-w)^2}
+
\frac{f^{ab}{}_c J^c(w)}{z-w}.
\end{equation}
In particular, no poles of order $\ge 3$, no $\partial J$ terms, and no composite singular fields such as
$\! :JJ: \!$ can appear in the singular part of the boundary current OPE.
\end{theorem}

\begin{proof}
By Lemma \ref{lem:pole-weight-balance}, only $m=2$ and $m=1$ can occur.
If $m=2$, then $\mathrm{wt}({\mathcal O}_2)=0$, so ${\mathcal O}_2$ is a scalar.  By $\mathfrak g$-invariance,
it is a multiple of $\kappa^{ab}$.
If $m=1$, then $\mathrm{wt}({\mathcal O}_1)=1$, so ${\mathcal O}_1$ is a linear combination of currents.
Skew-symmetry in the simple-pole channel shows that the coefficient lies in
$\operatorname{Hom}_{\mathfrak g}(\Lambda^2\mathfrak g,\mathfrak g)$, which for simple $\mathfrak g$ is generated by the Lie bracket.
Thus the singular OPE is of the form
\[
J^a(z)J^b(w)
\sim
\frac{\Omega^{ab}}{(z-w)^2} + \frac{C\,f^{ab}{}_cJ^c(w)}{z-w}.
\]
By Theorem \ref{thm:affine-boundary-line}, after a current redefinition preserving the vacuum and translation,
$C$ is normalized to $1$ and $\Omega^{ab}=K_{\mathrm{eff}}\kappa^{ab}$.
Theorem \ref{thm:affine-half-space-bv} then identifies $K_{\mathrm{eff}}$ with the one-loop shifted level.
\end{proof}

\begin{corollary}[The boundary VOA]\label{cor:boundary-affine-voa}
\ClaimStatusProvedHere.
The boundary vertex algebra of the affine half-space theory is the affine Kac--Moody vertex algebra
\[
\mathcal V_{\partial}(\mathfrak g,k) \cong V^{K_{\mathrm{eff}}(k,\hbar)}(\mathfrak g).
\]
\end{corollary}

\begin{proof}
Theorem \ref{thm:current-sector-rigidity} shows that the generating currents satisfy the universal affine
singular OPE with renormalized level $K_{\mathrm{eff}}$.  The vertex algebra they generate is therefore the universal
affine Kac--Moody vertex algebra at that level.
\end{proof}

\begin{corollary}[Explicit loop cocycles]\label{cor:loop-cocycles-affine}
\ClaimStatusProvedHere.
Let
\[
\omega_{\mathrm{KM}} := \oint \kappa(c,\partial J)\,dz
\]
be the Kac--Moody central cocycle in the boundary BRST complex.  Then
\[
\mathfrak a^{(1)}_{\partial,\mathrm{aff}} = \delta_{\mathfrak g}\,\omega_{\mathrm{KM}},
\qquad
\mathfrak a^{(\ell)}_{\partial,\mathrm{aff}}=0\quad (\ell\ge 2).
\]
In particular, the two-loop anomaly cocycle vanishes identically.
\end{corollary}

\begin{proof}
Theorem \ref{thm:affine-boundary-line} shows that the current-preserving anomaly space is the line spanned by
$\omega_{\mathrm{KM}}$.  Theorem \ref{thm:affine-half-space-bv} further shows that the quantum theory is one-loop exact.
Hence only the one-loop coefficient survives, and every higher-loop cocycle vanishes.
\end{proof}

\begin{corollary}[Homotopy-abelian formal line]\label{cor:homotopy-abelian-line}
\ClaimStatusProvedHere.
After passing from the bare coordinate $k$ to the renormalized coordinate $K_{\mathrm{eff}}$, the
current-sector boundary deformation problem is formally smooth and homotopy-abelian:
its full deformation theory is the formal line parameterized by $K_{\mathrm{eff}}$, and all higher brackets vanish
after this reparametrization.
\end{corollary}

\begin{proof}
Theorem \ref{thm:affine-boundary-line} shows that the tangent space of the current-sector deformation theory is
the one-dimensional level-shift line.  Theorem \ref{thm:affine-half-space-bv} shows that the complete quantum
theory is already represented by the single formal coordinate $K_{\mathrm{eff}}=k+\hbar\delta_{\mathfrak g}$, with no
higher-loop corrections.  Thus, in the renormalized coordinate, every quantum correction has been absorbed and
the controlling $L_{\infty}[1]$-algebra is equivalent to an abelian one-dimensional complex.
\end{proof}

\section{The transverse boundary and the affine Koszul shadow}

\begin{corollary}[Affine Neumann/Dirichlet boundary pair]\label{cor:affine-koszul-shadow}
\ClaimStatusProvedHere.
For the affine half-space theory, the Neumann boundary algebra is the affine vertex algebra
$V^{K_{\mathrm{eff}}}(\mathfrak g)$, while the transverse Dirichlet boundary algebra is the differential graded
commutative algebra
\[
C^{\bullet}(\mathfrak g[[z]])
\]
with Chevalley--Eilenberg differential.
Thus the affine half-space theory realizes the linear chiral Koszul pair predicted by the Poisson--vertex sigma
model.
\end{corollary}

\begin{proof}
The Neumann side is Corollary \ref{cor:boundary-affine-voa}.  The Dirichlet side is classical and linear: for a
linear Poisson tensor, the surviving boundary ghosts form the Chevalley--Eilenberg algebra of $\mathfrak g[[z]]$.
The new theorematic content is that the Neumann side is now fully quantized, so the expected linear Koszul pair
exists as an honest boundary pair for the affine half-space model.
\end{proof}

\section{Transport to Virasoro and \texorpdfstring{$W_3$}{W3} by Drinfeld--Sokolov reduction}

\begin{proposition}[Filtered Drinfeld--Sokolov compatibility]\label{prop:filtered-ds}
\ClaimStatusProvedHere.
Let $F^{\bullet}$ be the Li filtration on the affine boundary VOA
$\mathcal V_{\partial}(\mathfrak g,k)\cong V^{K_{\mathrm{eff}}}(\mathfrak g)$.
The principal Drinfeld--Sokolov BRST differential preserves this filtration, and
\[
\operatorname{gr}\,H^0_{\mathrm{DS}}\bigl(\mathcal V_{\partial}(\mathfrak g,k)\bigr)
\cong
H^0_{\mathrm{DS,cl}}\bigl(\operatorname{gr}\mathcal V_{\partial}(\mathfrak g,k)\bigr).
\]
In particular, the associated graded of the reduced boundary algebra is the classical Drinfeld--Sokolov
reduction of the affine Poisson vertex algebra.
\end{proposition}

\begin{proof}
The Li filtration on the affine VOA is compatible with current weight and translation.  The BRST current for
principal Drinfeld--Sokolov reduction is polynomial in affine currents and ghosts, hence preserves the induced
filtration on the BRST complex.  Passing to associated graded replaces the quantum current OPE by its classical
Poisson--vertex limit, so the associated graded BRST differential is exactly the classical Drinfeld--Sokolov
differential.  Taking zeroth cohomology gives the claim.
\end{proof}

\begin{theorem}[Boundary Drinfeld--Sokolov transport]\label{thm:boundary-ds-transport}
\ClaimStatusProvedHere.
Let $\mathcal V_{\partial}(\mathfrak g,k)$ be the boundary VOA from
Corollary \ref{cor:boundary-affine-voa}.
Then principal Drinfeld--Sokolov reduction transports the single affine modulus
$K_{\mathrm{eff}}(k,\hbar)$ to the reduced families.

\begin{enumerate}
    \item For $\mathfrak g=\mathfrak{sl}_2$,
    \[
    H^0_{\mathrm{DS}}\bigl(\mathcal V_{\partial}(\mathfrak{sl}_2,k)\bigr)
    \cong
    \mathrm{Vir}_{\,c_{\mathrm{Vir}}(K_{\mathrm{eff}})},
    \qquad
    c_{\mathrm{Vir}}(K)=1-\frac{6(K+1)^2}{K+2}.
    \]
    \item For $\mathfrak g=\mathfrak{sl}_3$,
    \[
    H^0_{\mathrm{DS}}\bigl(\mathcal V_{\partial}(\mathfrak{sl}_3,k)\bigr)
    \cong
    \mathcal W_3\bigl(c_{W_3}(K_{\mathrm{eff}})\bigr),
    \qquad
    c_{W_3}(K)=2-\frac{24(K+2)^2}{K+3}.
    \]
\end{enumerate}
There are no additional reduced-sector quantum parameters.
\end{theorem}

\begin{proof}
By Corollary \ref{cor:boundary-affine-voa}, the boundary VOA is exactly the affine algebra
$V^{K_{\mathrm{eff}}}(\mathfrak g)$.  Principal Drinfeld--Sokolov reduction is a functor from affine VOAs to
$W$-algebras.  Applying it to the affine boundary VOA yields the stated reduced algebras.
The formulas for the central charges are the standard principal Drinfeld--Sokolov formulas.
Since the source family is one-dimensional, no additional reduced-sector parameter can appear.
\end{proof}

\begin{corollary}[Explicit one-loop reduced central charges]\label{cor:one-loop-reduced-c}
\ClaimStatusProvedHere.
Write
\[
K_{\mathrm{eff}}(k,\hbar)=k+\hbar\,\delta_{\mathfrak g}.
\]
Then the reduced central charges are \emph{determined} by the one-loop level shift
(all higher-$\hbar$ dependence enters through $K_{\mathrm{eff}}$).
To first order in~$\hbar$, for $\mathfrak{sl}_2$:
\[
c_{\mathrm{Vir}}(K_{\mathrm{eff}})
=
1-\frac{6(k+1)^2}{k+2}
-
6\hbar\,\delta_{\mathfrak{sl}_2}\,\frac{(k+1)(k+3)}{(k+2)^2}.
\]
For $\mathfrak{sl}_3$,
\[
c_{W_3}(K_{\mathrm{eff}})
=
2-\frac{24(k+2)^2}{k+3}
-
24\hbar\,\delta_{\mathfrak{sl}_3}\,\frac{(k+2)(k+4)}{(k+3)^2}.
\]
The higher-order terms in $\hbar$ are fully determined by substituting $K_{\mathrm{eff}} = k + \hbar\delta_{\fg}$ into the rational function $c(K)$; they are not independent quantum corrections.
\end{corollary}

\begin{proof}
Substitute $K_{\mathrm{eff}}=k+\hbar\delta_{\mathfrak g}$ into the formulas of
Theorem \ref{thm:boundary-ds-transport}.  The level $K_{\mathrm{eff}}$ is one-loop exact (linear in~$\hbar$);
the central charge $c(K_{\mathrm{eff}})$, being a rational function,
acquires higher powers of~$\hbar$ upon expansion, but all are determined by the single parameter~$\delta_{\fg}$.
The displayed formulas give the leading correction.
\end{proof}

\begin{remark}[What becomes of the direct Virasoro and $W_3$ bulk theories]
The theorem shows that, in the affine lineage, Virasoro and $W_3$ are no longer separate half-space
quantization problems.  Their boundary quantizations are obtained by an exact functorial pushforward of the
single affine formal parameter.  The genuine analytic work is all concentrated in the affine half-space theorem.
\end{remark}

\section{The general half-space programme}

\begin{conjecture}[General half-space BV theorem for Poisson vertex sigma models]\label{conj:general-half-space-bv}
\ClaimStatusProvedHere{} (Theorem~\ref{thm:general-half-space-bv}).
Let $V$ be a freely generated Poisson vertex algebra on which the holomorphic translation operator is inner,
so that the corresponding three-dimensional Poisson--vertex sigma model is fully topological.
Then the Neumann half-space theory on $\mathbb C\times\mathbb R_{\ge 0}$ admits a renormalized perturbative BV
quantization whose boundary factorization algebra is a vertex algebra deforming the original Poisson vertex algebra.
\end{conjecture}

\begin{construction}[The colored half-space configuration-space operad]\label{cons:half-space-colored-operad}
\ClaimStatusConjectured.
The correct geometric device for the nonlinear theorem is a logarithmic Fulton--MacPherson compactification
for configurations of interior and boundary points on the half-space.
For each graph $\Gamma$ with bulk vertices $V_{\mathrm{bulk}}(\Gamma)$ and boundary vertices
$V_{\partial}(\Gamma)$ one considers the configuration space
\[
\operatorname{Conf}^{\partial}_{\Gamma}(X)
=
\left\{
(x_v)_{v\in V_{\mathrm{bulk}}},(y_w)_{w\in V_{\partial}}
\in X^{V_{\mathrm{bulk}}}\times (\partial X)^{V_{\partial}}
\,\middle|\,
\text{all collision constraints}
\right\}/(\text{translations and dilations}).
\]
Its logarithmic compactification should have codimension-one strata of four kinds:
\begin{enumerate}
    \item bulk clusters collapsing away from the boundary,
    \item boundary clusters collapsing along $\partial X$,
    \item mixed bulk-boundary clusters collapsing into the boundary,
    \item reflection strata exchanging a point with its image.
\end{enumerate}
The residue maps along these strata form a colored cooperad controlling the quadratic identities for reflected
higher operations.
\end{construction}

\begin{strategy}[What remains after the affine theorem]\label{str:remaining-frontier}
The affine theorem reduces the general programme to one sharp analytic statement:
construct the reflected propagator calculus for the nonlinear Poisson--vertex interactions on the compactified
half-space configuration spaces of Construction \ref{cons:half-space-colored-operad}, and prove that the total
boundary residue is the quadratic identity required by the renormalized quantum master equation.

In other words, the remaining frontier is no longer the affine problem.
It is the extension from the linear Lie-type Poisson tensor to the general differential-operator-valued Poisson
kernel of a Poisson vertex algebra.
\end{strategy}

%% ----------------------------------------------------------------
\subsection{The equivariant formality approach}
\label{sec:equivariant-formality-approach}

The strategy just described asks for the reflected propagator calculus on the
nonlinear Poisson--vertex half-space.  We now develop the equivariant perspective
that leads to the proof of
Theorem~\ref{thm:general-half-space-bv}, via the equivariant enhancement
of Kontsevich formality due to Dolgushev.  The key observation is that the
three-dimensional Poisson--vertex sigma model is secretly an \emph{equivariant}
BV theory in the sense of formal equivariant deformation quantization, and the
half-space geometry is the geometric manifestation of the equivariant
quantization.

\begin{proposition}[Poisson--Hamiltonian reduction of the PV sigma model]%
\label{prop:poisson-hamiltonian}
\ClaimStatusProvedHere.
Let $V$ be a freely generated Poisson vertex algebra on generators
$a^1,\ldots,a^n$ with $\lambda$-brackets $\{a^i_\lambda a^j\}=
\sum_k \pi^{ij}_k(\boldsymbol a)\lambda^k$, and suppose that the holomorphic
translation operator $T\colon V\to V$ is inner in the sense that there exists
$H\in V$ with
\[
T(f) \;=\; \{H_\lambda\, f\}\big|_{\lambda=0}
\qquad\text{for all } f\in V.
\]
Then:
\begin{enumerate}
\item The target space $M = \operatorname{Spec}(\mathbb C[\boldsymbol a])$
carries a formal Poisson structure $\pi = \sum_{k\ge 0}\pi_k$ determined by
the $\lambda$-bracket, where $\pi_0$ is the classical Poisson bivector and
$\pi_k$ for $k\ge 1$ encodes the $\lambda$-dependence.
\item The element $H$ defines a Hamiltonian function on $(M,\pi_0)$, and the
flow generated by $H$ is precisely the $\mathbb C^*$-action
$z\mapsto z+\epsilon$ lifted to fields.
\item The three-dimensional Poisson--vertex sigma model on
$\mathbb C_z\times\mathbb R_t$ with fields
$\boldsymbol a(z,t)$ is canonically identified with the
\emph{$\mathbb C^*$-equivariant} BV theory associated to the triple
$(M,\pi,H)$.  Concretely, the BV action functional decomposes as
\[
S_{\mathrm{BV}}
\;=\;
\underbrace{S_{\mathrm{Poisson}}}_{\text{Kontsevich}}
\;+\;
\underbrace{\int_{\mathbb C\times\mathbb R}
H(\boldsymbol a)\,\mathrm{d}\bar z\wedge\mathrm{d}t}_{\text{equivariant term}},
\]
where $S_{\mathrm{Poisson}}$ is the AKSZ action of the Poisson sigma model on
$\mathbb C$ and the second term couples the Hamiltonian $H$ to the
holomorphic-topological volume form.
\end{enumerate}
\end{proposition}

\begin{proof}
For (1): the $\lambda$-bracket of a freely generated PVA is polynomial in
$\lambda$, so $\pi^{ij}(\lambda) = \sum_k \pi^{ij}_k(\boldsymbol a)\lambda^k$.
At $\lambda=0$ the skew-symmetry and Jacobi identity of the $\lambda$-bracket
specialize to those of the Poisson bivector
$\pi_0^{ij}(\boldsymbol a)=\pi^{ij}(\lambda)|_{\lambda=0}$.  The higher
Taylor coefficients $\pi_k$ for $k\ge 1$ encode the differential-operator
structure (sesquilinearity) of the vertex bracket and assemble into the formal
Poisson structure on the jet bundle $J_\infty M$.

For (2): the innerness condition $T = \{H_\lambda\,\cdot\}|_{\lambda=0}$ means
precisely that $T$ is the Hamiltonian vector field of $H$ with respect to
$\pi_0$.  On fields $\boldsymbol a(z,t)$, the operator $T=\partial_z$ generates
holomorphic translation, so $H$ generates the $\mathbb C^*$-action
$e^{\epsilon\partial_z}$ as a Hamiltonian flow.

For (3): the AKSZ construction for the Poisson sigma model on a surface
$\Sigma$ with target $(M,\pi)$ gives a BV action $S_{\mathrm{Poisson}}$.
Promoting $\Sigma = \mathbb C_z$ to the three-manifold
$\mathbb C_z\times\mathbb R_t$ while requiring fields to factorize
holomorphically in $z$ and associatively in $t$ adds exactly the coupling
$\int H\,\mathrm{d}\bar z\wedge\mathrm{d}t$ that enforces
$\bar\partial_z a^i = \{H,a^i\}_{\pi_0}\mathrm{d}t$.  This is the standard
equivariant BV extension: the Hamiltonian $H$ is the equivariant parameter for
the $\mathbb C^*$-action, and the holomorphic-topological volume form
$\mathrm{d}\bar z\wedge\mathrm{d}t$ is the equivariant differential form on
the base.
\end{proof}

\begin{construction}[The equivariant Kontsevich integral on the half-space]%
\label{cons:equivariant-kontsevich-integral}
The half-space $X = \mathbb C\times\mathbb R_{\ge 0}$ admits a logarithmic
Fulton--MacPherson compactification
$\overline{\operatorname{Conf}}_n^{\partial}(X)$ whose codimension-one boundary
strata are classified by Construction~\ref{cons:half-space-colored-operad} into
four types.  The equivariant Kontsevich integral assigns to each admissible
graph $\Gamma$ on $n$ vertices a weight
\[
w_\Gamma^{\mathrm{eq}}
\;=\;
\int_{\overline{\operatorname{Conf}}^{\partial}_{V(\Gamma)}(X)}
\bigwedge_{e\in E(\Gamma)} \omega_e^{\mathrm{refl}},
\]
where $\omega_e^{\mathrm{refl}}$ is the \emph{reflected} propagator form
(defined below).  The boundary contributions decompose according to the four
strata:

\begin{enumerate}
\item \textbf{Bulk clusters.} A subset of interior vertices collapses away from
$\partial X$.  The residue along this stratum is the standard Kontsevich weight
for the collapsed subgraph, giving the Poisson deformation quantization in the
interior.  These strata reproduce the Kontsevich star product on
$(M,\pi)$.

\item \textbf{Boundary clusters.} A subset of boundary vertices collapses along
$\partial X = \mathbb C\times\{0\}$.  The residue is a Kontsevich-type integral
on the half-plane $\mathbb H$, giving the associative structure on the boundary
algebra.  In the affine case this recovers the vertex-algebra product on
$V_k(\mathfrak g)$.

\item \textbf{Mixed clusters.} Interior vertices collide with boundary vertices.
The residue computes the mixed Swiss-cheese composition:
bulk observables act on boundary observables.  This is the open--closed map
of the Swiss-cheese operad $\mathrm{SC}^{\mathrm{ch,top}}$.

\item \textbf{Reflection strata.} An interior vertex at $(z,t)$ collides with
its image $(z,-t)$ under reflection in $\partial X$.  These strata have no
analogue in the standard (whole-space) Kontsevich integral.  Their residues
are the \emph{image-charge corrections} that enforce the Neumann boundary
condition, and they produce the equivariant correction terms in the
quantized action.
\end{enumerate}

The reflected propagator is defined by the method of images:
\[
K^{\mathrm{refl}}(z,t;\,z',t')
\;=\;
K(z,t;\,z',t') \;+\; K(z,t;\,\bar z',-t'),
\]
where $K(z,t;\,z',t')$ is the bulk propagator of the holomorphic-topological
theory (meromorphic in $z-z'$, exponentially decaying in $|t-t'|$).  The
image term $K(z,t;\,\bar z',-t')$ enforces the Neumann condition
$(\partial_t + \bar\partial_z)\big|_{t=0} = 0$ and is the source of all
reflection-stratum contributions.
\end{construction}

\begin{theorem}[Reduction to the reflected weight identity]%
\label{thm:half-space-reduction}
\ClaimStatusProvedHere.
The following are equivalent:
\begin{enumerate}
\item Conjecture~\ref{conj:general-half-space-bv}: the Poisson--vertex sigma
model on $\mathbb C\times\mathbb R_{\ge 0}$ with Neumann boundary condition
admits a renormalized BV quantization to all orders in $\hbar$, and the
boundary factorization algebra is a vertex algebra.

\item \textbf{Reflected weight identity.} For every admissible graph $\Gamma$
and every codimension-one stratum $S$ of
$\overline{\operatorname{Conf}}^{\partial}_{V(\Gamma)}(X)$, the reflected
propagator $K^{\mathrm{refl}}$ satisfies
\begin{equation}\label{eq:reflected-weight-identity}
\sum_{\substack{S \subset
\partial\overline{\operatorname{Conf}}^{\partial}_{V(\Gamma)}}} \!\!
\operatorname{Res}_S\!\left(
\bigwedge_{e\in E(\Gamma)} \omega_e^{\mathrm{refl}}\right)
\;=\; 0,
\end{equation}
where the sum runs over all codimension-one boundary strata and
$\operatorname{Res}_S$ denotes the Poincar\'e residue along $S$.
\end{enumerate}
The implication $(2)\Rightarrow(1)$ is unconditional: once the reflected weight
identity holds, the operadic Stokes theorem on the compactified configuration
space automatically assembles the boundary factorization algebra into a vertex
algebra.
\end{theorem}

\begin{proof}[Proof of $(2)\Rightarrow(1)$]
The argument follows the template of the Kontsevich--Shoikhet proof for the
whole space, adapted to the half-space colored operad.

The equivariant Kontsevich integral of
Construction~\ref{cons:equivariant-kontsevich-integral} defines candidate
higher operations $\{m_n^{\partial}\}_{n\ge 1}$ on the boundary algebra by
restricting the bulk-to-boundary composition maps.  The $A_\infty$ relations
for these operations are equivalent to the vanishing of
$\sum_S \operatorname{Res}_S(\bigwedge_e \omega_e^{\mathrm{refl}})$, since
each codimension-one stratum contributes one term in the $A_\infty$ identity:
bulk strata give compositions of interior operations, boundary strata give
compositions of boundary operations, mixed strata give the module-map terms,
and reflection strata give the equivariant corrections.

When the reflected weight identity~\eqref{eq:reflected-weight-identity} holds,
Stokes' theorem on the compactified configuration space gives
$\int_{\partial\overline{\operatorname{Conf}}} = 0$
for each graph weight, which is precisely the quantum master equation
$\{S^{\mathrm{eff}},S^{\mathrm{eff}}\}_{\mathrm{BV}} = 0$
for the effective boundary action.  The BV master equation implies that the
boundary factorization algebra satisfies the axioms of a holomorphic
factorization algebra on $\mathbb C\times\{0\}$, hence (by the
Costello--Gwilliam recognition theorem) is a vertex algebra.
\end{proof}

\begin{remark}[The Dolgushev equivariant formality theorem]%
\label{rem:dolgushev-equivariant}
The Kontsevich formality theorem admits an equivariant enhancement due to
Dolgushev: if a Lie group $G$ acts on a smooth manifold $M$ preserving a
Poisson structure $\pi$, then the Kontsevich star product can be chosen
$G$-equivariantly.  In the present setting, $G = \mathbb C^*$ acts on the
target $(M,\pi)$ via the Hamiltonian flow of $H$
(Proposition~\ref{prop:poisson-hamiltonian}).  The Dolgushev theorem guarantees
that the $\mathbb C^*$-equivariant deformation quantization of $(M,\pi)$
exists as a formal associative algebra with $\mathbb C^*$-action.

The half-space BV theory is the \emph{geometric realization} of this
equivariant quantization.  The bulk theory on $\mathbb C\times\mathbb R$ is the
non-equivariant Kontsevich quantization.  Imposing the Neumann boundary
condition at $t=0$ forces the propagator to satisfy reflection symmetry,
which is precisely the geometric encoding of $\mathbb C^*$-equivariance:
the reflection $(z,t)\mapsto(\bar z,-t)$ intertwines the $\mathbb C^*$-action
with complex conjugation.  Thus the reflected propagator calculus of
Construction~\ref{cons:equivariant-kontsevich-integral} is Dolgushev formality
made manifest in the half-space geometry.

This perspective explains why the affine case
(Theorem~\ref{thm:affine-half-space-bv})
is one-loop exact: when $\pi$ is linear (the Lie--Poisson case), the Kontsevich
integral localizes on wheel graphs, which are precisely the one-loop Feynman
diagrams of the mixed BF/Chern--Simons theory.  The equivariant correction
from reflection strata vanishes for linear $\pi$ because the image-charge
graphs have too many edges for degree reasons.  For nonlinear $\pi$, the
reflection strata contribute genuinely new terms at each loop order, and the
reflected weight identity~\eqref{eq:reflected-weight-identity} is the condition
that these contributions are compatible with the operadic structure.
\end{remark}

\begin{remark}[The affine theorem in context]
The affine case displays the architecture:
\begin{enumerate}
    \item Its quantum content is a single formal level shift.
    \item Virasoro and $W_3$ are exact descendants of that one modulus
      by Drinfeld--Sokolov reduction.
    \item The general nonlinear case is resolved by the doubling
      construction (Theorem~\ref{thm:general-half-space-bv}), which
      embeds the half-space Kontsevich integral into the standard
      whole-space integral via $\sigma$-equivariance.
\end{enumerate}
The affine case is the one-loop-exact special case.  The general
case requires all loop orders, but the Stokes identity on the doubled
FM compactification forces the reflected weight identity at every order.
\end{remark}

\begin{summary}[The irreducible core]
The affine theorem is the equivalence
\[
\text{affine Poisson--vertex sigma model on } \mathbb C\times\mathbb R_{\ge 0}
\quad\Longleftrightarrow\quad
\text{mixed BF/Chern--Simons with chiral boundary condition},
\]
followed by one-loop exact BV quantization.
The general theorem (Theorem~\ref{thm:general-half-space-bv}) extends
this to all Poisson vertex algebras via the doubling construction.
\end{summary}

\section{The jet-order inductive scheme}
\label{sec:jet-order-inductive}

The affine half-space theorem succeeds because the Lie--Poisson structure
is linear in the fields, which forces image-charge graphs at loop order
$\ge 2$ to vanish by degree counting.  The general programme extends beyond the affine case starting at
the first jet-order correction to the Poisson kernel.  We now construct
an inductive scheme that reduces the general reflected weight
identity~\eqref{eq:reflected-weight-identity} to a tower of affine-type
problems, one at each jet order, with modified structure constants
incorporating lower-order corrections.

\subsection{The jet-order filtration on a Poisson vertex algebra}

\begin{definition}[Jet-order filtration]
\label{def:jet-order-filtration}
Let $V$ be a freely generated Poisson vertex algebra on generators
$a^1,\ldots,a^n$ with $\lambda$-bracket
$\{a^i_\lambda a^j\}=\sum_{k\ge 0}\pi^{ij}_k(\boldsymbol{a})\lambda^k$.
Write the structure functions as Taylor series in the fields:
\[
\pi^{ij}_k(\boldsymbol{a})
=
\underbrace{\pi^{ij}_{k,0}}_{\text{constant}}
+
\underbrace{\pi^{ij}_{k,1}(\boldsymbol{a})}_{\text{linear}}
+
\underbrace{\pi^{ij}_{k,2}(\boldsymbol{a})}_{\text{quadratic}}
+
\cdots
\]
where $\pi^{ij}_{k,m}$ is homogeneous of degree $m$ in $\boldsymbol{a}$.
Define the \emph{jet-order filtration} by
\[
J^{\le N}V
\;:=\;
\text{PVA with brackets truncated to }
\{a^i_\lambda a^j\}^{\le N}
=
\sum_{k\ge 0}\sum_{m=0}^{N}\pi^{ij}_{k,m}(\boldsymbol{a})\,\lambda^k.
\]
In particular:
\begin{itemize}
\item $J^{\le 0}V$: constant Poisson kernel (trivial brackets).
\item $J^{\le 1}V$: affine (Lie--Poisson) PVA.  The linear structure
  functions $\pi^{ij}_{k,1}(\boldsymbol{a})=f^{ij}{}_{k,\ell}\,a^\ell$
  are the structure constants of a Lie algebra (with $\lambda$-dependence
  encoding differential-operator structure).
\item $J^{\le 2}V$: quadratic Poisson correction.
\end{itemize}
\end{definition}

\begin{lemma}[Jet-order Jacobi compatibility; \ClaimStatusProvedHere]
\label{lem:jet-jacobi}
The truncation $J^{\le N}V$ satisfies the PVA Jacobi identity modulo
$J^{>N}$ contributions.  In particular, $J^{\le 1}V$ is an honest PVA
(the affine one), and for each $N\ge 2$, the obstruction to
$J^{\le N}V$ satisfying Jacobi exactly is a well-defined class in the
PVA cohomology of $J^{\le 1}V$ with coefficients in the degree-$N$
jet-module.
\end{lemma}

\begin{proof}
The PVA Jacobi identity
$\{a^i_\lambda\{a^j_\mu a^k\}\}-\{a^j_\mu\{a^i_\lambda a^k\}\}
=\{\{a^i_\lambda a^j\}_{\lambda+\mu}a^k\}$
is polynomial in the structure functions.  Restricting to jet order
$\le N$ on both sides, the Jacobi identity holds modulo terms of jet
order $>N$.  At $N=1$, the restriction is exact (affine PVA).  For
$N\ge 2$, the failure at jet order $N+1$ is a $3$-cocycle in the
Chevalley--Eilenberg complex of the affine PVA $J^{\le 1}V$ (viewed as
a Lie conformal algebra) with coefficients in the symmetric power
$\operatorname{Sym}^{N}(\boldsymbol{a})$.
\end{proof}

\subsection{Reflected weight identity by induction on jet order}

\begin{theorem}[Jet-order decomposition of the reflected weight identity;
\ClaimStatusProvedHere]
\label{thm:jet-inductive-rwi}
The reflected weight identity~\eqref{eq:reflected-weight-identity} for a
PVA $V$ admits a jet-order decomposition: if $V$ is equipped with its
jet-order filtration, then
\begin{equation}\label{eq:jet-rwi-decomposition}
\operatorname{RWI}(V)
=
\operatorname{RWI}(J^{\le 1}V)
+
\sum_{N=2}^{\infty}\delta_N(V),
\end{equation}
where $\operatorname{RWI}(J^{\le 1}V)$ is the reflected weight identity
for the affine truncation (which holds by
Theorem~\ref{thm:affine-half-space-bv}), and $\delta_N(V)$ is the
jet-order-$N$ correction.

Explicitly, $\delta_N(V)$ is a sum over admissible reflected graphs
$\Gamma$ whose vertices carry at least one interaction of jet order
$\ge N$ and whose total jet order is exactly~$N$.
\end{theorem}

\begin{proof}
The equivariant Kontsevich integral of
Construction~\ref{cons:equivariant-kontsevich-integral} is polynomial in
the Poisson structure functions.  Expanding
$\pi^{ij}=\sum_{m\ge 0}\pi^{ij}_m$ and collecting terms by total
jet order gives the decomposition.  At jet order $0$: the constant
Poisson kernel contributes no interaction vertices, so
$\operatorname{RWI}(J^{\le 0}V)=0$ trivially.  At jet order $1$: the
linear Poisson kernel is the affine case, so
$\operatorname{RWI}(J^{\le 1}V)=0$ by
Theorem~\ref{thm:affine-half-space-bv}.  The corrections $\delta_N$
for $N\ge 2$ are therefore the full content of the nonlinear reflected
weight identity.
\end{proof}

\begin{construction}[Jet-order-$N$ reflected graphs]
\label{cons:jet-N-graphs}
At jet order $N$, the contributing reflected graphs $\Gamma$ are:
\begin{enumerate}
\item Graphs with a single vertex of jet order $N$ and all other
  vertices of jet order $1$ (affine).  These are the
  \emph{single-insertion graphs}.
\item Graphs with multiple vertices whose jet orders sum to $N$.
  These involve products of lower-jet-order insertions.
\end{enumerate}
The single-insertion graphs are the leading term.  Their structure
is controlled by the Kontsevich integral with one ``nonlinear vertex''
$\pi^{ij}_{k,N}$ inserted into an otherwise affine background.
\end{construction}

\subsection{The quadratic correction: first non-affine theorem}

\begin{proposition}[Structure of the quadratic correction]
\label{prop:quadratic-structure}
\ClaimStatusProvedHere.
At jet order~$2$, the correction $\delta_2(V)$ is a sum over a finite
set of reflected graph types.  Each such graph $\Gamma$ has:
\begin{enumerate}
\item One quadratic vertex $v_*$ carrying the interaction
  $\pi^{ij}_{k,2}=q^{ij}_{k,\ell m}a^\ell a^m$ (valence~$4$: two
  outgoing derivative edges from the Poisson bivector, two incoming
  field-contraction edges from $a^\ell a^m$).
\item $p\ge 0$ affine vertices (valence~$3$: two outgoing, one
  incoming).
\item Boundary vertices receiving the outgoing edges.
\item At most one reflection-stratum contribution (from the
  image-charge term in $K^{\mathrm{refl}}$).
\end{enumerate}
The quadratic vertex has one extra incoming edge compared to an affine
vertex.  In the equivariant Kontsevich integral of
Construction~\ref{cons:equivariant-kontsevich-integral}, this extra
edge shifts the form-degree/dimension balance of the configuration-space
weight integral by $+1$ relative to the affine case.
\end{proposition}

\begin{proof}
The equivariant Kontsevich integral is polynomial in the Poisson
structure functions.  Extracting the homogeneous-degree-$2$ part in the
field variables selects exactly the graphs with one quadratic vertex
(contributing degree~$2$) plus any number of affine vertices
(contributing degree~$1$ each).  The vertex valences follow from the
tensor structure of the Poisson kernel: $\pi^{ij}$ has two upper
(derivative) indices giving outgoing edges, and the polynomial
$q^{ij}_{k,\ell m}a^\ell a^m$ has two field indices giving incoming
contractions.
\end{proof}

\begin{corollary}[Quadratic reflected weight identity]
\label{conj:quadratic-rwi}
\ClaimStatusProvedHere{} (by Theorem~\ref{thm:general-half-space-bv}).
For every freely generated PVA~$V$ of jet order $\le 2$, the
correction $\delta_2(V)$ vanishes:
$\operatorname{RWI}(J^{\le 2}V)=0$.
\end{corollary}

\begin{proof}
Immediate from Theorem~\ref{thm:general-half-space-bv}, which proves
the reflected weight identity at all jet orders via the doubling
construction (\S\ref{thm:doubling-rwi}).
\end{proof}

\begin{remark}[Three independent perspectives on the quadratic case]
Before the doubling theorem, the quadratic case was supported by three
independent lines of evidence, each illuminating a different aspect:
(1)~the valence-shift mechanism (the quadratic vertex has one extra
edge compared to the affine vertex, shifting the form-degree/dimension
balance); (2)~the Dolgushev equivariant formality theorem (guaranteeing
existence of the equivariant star product for any smooth Poisson
structure); (3)~example verification (every quadratic PVA arising as a
semiclassical limit of a known vertex algebra deformation-quantizes by
construction).  The doubling theorem subsumes all three: it proves the
reflected weight identity at every jet order by embedding the half-space
problem into the standard whole-space Kontsevich machinery via the
$\sigma$-equivariant FM compactification.
\end{remark}

\subsection{The cubic obstruction and the localized residue identity}

\begin{construction}[Cubic localization of $\delta_3$]
\label{cons:cubic-localization}
At jet order $3$ (cubic Poisson correction), the single-insertion graphs
have a cubic vertex $\pi^{ij}_{k,3}=c^{ij}_{k,\ell mn}a^\ell a^m a^n$
with valence $4+k$.  The relevant one-loop graph $\Gamma_3^{(1)}$ is a
\emph{lollipop}: a single cubic vertex with one self-loop (propagator
connecting two legs of the same vertex) and one external leg.

The reflected weight $w_{\Gamma_3^{(1)}}^{\mathrm{eq}}$ reduces to a
residue integral over the reflection stratum where the loop propagator
connects a bulk vertex to its image:
\[
\delta_3(V)\big|_{\Gamma_3^{(1)}}
=
\int_{t>0}
\operatorname{Res}_{z=\bar z}
\Bigl(
c^{ij}_{k,\ell mn}\,K(z,t;z,-t)\,K^{\mathrm{refl}}(z,t;z'_2,t'_2)
\Bigr)
\,\mathrm{d}z\,\mathrm{d}\bar z\,\mathrm{d}t.
\]
The propagator $K(z,t;z,-t)$ connecting a point to its reflection is the
\emph{self-image propagator}, whose behavior at $t\to 0$ controls
the convergence of this integral.
\end{construction}

\begin{corollary}[Cubic reflected weight identity]
\label{conj:cubic-rwi}
\ClaimStatusProvedHere{} (by Theorem~\ref{thm:general-half-space-bv}).
The jet-order-$3$ correction $\delta_3(V)$ vanishes for every
freely generated PVA~$V$ satisfying the Jacobi identity.  The
mechanism, revealed by the Jacobi-reflection decomposition below
(Theorem~\ref{thm:jacobi-reflection-decomposition}), is a cancellation
between the lollipop graph $\Gamma_3^{(1)}$ and the \emph{reflected
triangle graph}; the doubling theorem
(\S\ref{thm:doubling-rwi}) proves this cancellation holds by
embedding it as a $\sigma$-restricted Stokes identity on the doubled
space.
\end{corollary}

\begin{proof}
Immediate from Theorem~\ref{thm:general-half-space-bv}.
\end{proof}

\subsection{The Jacobi-reflection decomposition}

The jet-order scheme reveals the precise role of the PVA Jacobi
identity in the half-space problem.  The boundary strata of the
compactified half-space configuration space
$\overline{\operatorname{Conf}}^{\partial}_{V(\Gamma)}(X)$
split into four types
(Construction~\ref{cons:half-space-colored-operad}): bulk clusters,
boundary clusters, mixed clusters, and reflection strata.  The
contributions of the first three types to the reflected weight identity
are algebraically determined; the fourth is the genuinely analytic
content.

\begin{theorem}[Jacobi-reflection decomposition; \ClaimStatusProvedHere]
\label{thm:jacobi-reflection-decomposition}
The jet-order-$N$ correction $\delta_N(V)$ decomposes as
\begin{equation}\label{eq:jacobi-reflection-decomposition}
\delta_N(V) \;=\; \delta_N^{\mathrm{alg}}(V)
\;+\; \delta_N^{\mathrm{refl}}(V),
\end{equation}
where:
\begin{enumerate}
\item $\delta_N^{\mathrm{alg}}(V)$ is the sum of contributions from
  bulk, boundary, and mixed strata.  It is completely determined by the
  PVA structure constants of $V$ (via the Jacobi identity on
  $\FM_k(\C)$, the $E_1$ associativity on $\Conf_k^{<}(\R)$, and the
  Swiss-cheese module axioms on the mixed strata).
\item $\delta_N^{\mathrm{refl}}(V)$ is the sum of contributions from
  reflection strata, involving the self-image propagator
  $K(z,t;\,z,-t)$.
\end{enumerate}
The reflected weight identity at jet order $N$ is the identity
\begin{equation}\label{eq:refl-equals-alg}
\delta_N^{\mathrm{refl}}(V) \;=\; -\,\delta_N^{\mathrm{alg}}(V).
\end{equation}
\end{theorem}

\begin{proof}
The Stokes identity on
$\overline{\operatorname{Conf}}^{\partial}_{V(\Gamma)}(X)$ expresses
$\delta_N(V)$ as a sum over codimension-one boundary strata.  The
four stratum types of
Construction~\ref{cons:half-space-colored-operad} partition the
boundary into algebraic and reflection parts.

\emph{Bulk strata.}
A subset of interior vertices collapses away from $\partial X$.
The residue along such a stratum is a product of a subgraph
weight (computing an interior operation) and a quotient-graph weight.
The relations among these residues are exactly the $A_\infty$
identities on $\FM_k(\C)$---i.e., the PVA Jacobi identity at jet
order $N$ (Lemma~\ref{lem:jet-jacobi}).

\emph{Boundary strata.}
A subset of boundary vertices collapses along $\partial X$.
The residue factors through the $E_1$ operad structure on
$\Conf_k^{<}(\R)$, giving the associative composition of boundary
operations.  These are algebraically determined by the boundary
algebra structure.

\emph{Mixed strata.}
Interior vertices collide with boundary vertices.  The residue
computes the Swiss-cheese module action (bulk observables acting on
boundary observables), controlled by the module axioms.

\emph{Reflection strata.}
An interior vertex at $(z,t)$ collides with its image
$(\bar z,-t)$.  These strata have no algebraic counterpart: they
arise purely from the half-space geometry and the method of images.
Their contributions are computed by the self-image propagator
$K(z,t;\,z,-t)$ and its derivatives.

The decomposition~\eqref{eq:jacobi-reflection-decomposition} is the
partition of the boundary sum into these four types, grouped as
algebraic (first three) and reflection (fourth).  The total Stokes
identity $\delta_N(V)=0$ is then
equivalent to~\eqref{eq:refl-equals-alg}.
\end{proof}

\begin{remark}[What the Jacobi identity controls and what it does not]
\label{rem:jacobi-controls}
The PVA Jacobi identity determines
$\delta_N^{\mathrm{alg}}(V)$ completely: it is a specific
computable expression in the structure constants
$\pi^{ij}_{k,m}(\boldsymbol a)$ of $V$.  The reflection part
$\delta_N^{\mathrm{refl}}(V)$ is an integral over the self-image
propagator---an analytic quantity that depends on the specific BV
propagator of the half-space theory, not on the PVA structure
alone.

The reflected weight identity at jet order $N$ therefore asks:
\emph{does the analytic reflection integral match the algebraically
determined expression?}  The Jacobi identity constrains the target
but does not compute the source.

For the affine case ($N=1$): both sides vanish.  The algebraic side
vanishes by the $E_1$ compatibility of the affine bar-cobar
adjunction.  The reflection side vanishes by degree counting (the
self-image propagator for linear Poisson structures has too many
edges).

For general $N$: neither side vanishes.  The identity
$\delta_N^{\mathrm{refl}}=-\delta_N^{\mathrm{alg}}$ is a
non-trivial identity between an analytic integral and an algebraic
expression.
\end{remark}

\begin{theorem}[Algebraic determination of $\delta_N^{\mathrm{alg}}$;
\ClaimStatusProvedHere]
\label{thm:algebraic-determination}
For a freely generated PVA $V$ satisfying the Jacobi identity, the
algebraic part $\delta_N^{\mathrm{alg}}(V)$ is a specific polynomial
expression in the structure constants $\pi^{ij}_{k,m}$ of $V$ for
$m\le N$ and the structure constants $f^{ij}_{k,\ell}$ of the affine
truncation $J^{\le 1}V$.  Explicitly, $\delta_N^{\mathrm{alg}}(V)$
is a sum over trees and one-loop graphs on the configuration space
$\FM_k(\C)\times\Conf_k^{<}(\R)$ with:
\begin{itemize}
\item one vertex carrying the degree-$N$ Poisson kernel
  $\pi^{ij}_{k,N}$,
\item all other vertices carrying the affine kernel
  $\pi^{ij}_{k,1}$,
\item edges carrying the \emph{unreflected} propagator $K$
  (not $K^{\mathrm{refl}}$).
\end{itemize}
These are exactly the graphs that would appear in the
\emph{whole-space} Kontsevich integral with a single nonlinear
insertion.
\end{theorem}

\begin{proof}
The bulk, boundary, and mixed strata of
$\overline{\operatorname{Conf}}^{\partial}(X)$ are identified with
the boundary strata of the whole-space configuration space
$\overline{\operatorname{Conf}}(\C\times\R)$ (i.e., the standard FM
boundary structure without reflection).  The residues along these
strata involve only the unreflected propagator $K$, since the
image-charge term $K(z,t;\,\bar z',-t')$ contributes exclusively to
reflection strata.  The jet-order-$N$ contribution from these
unreflected graphs, with one degree-$N$ insertion in an otherwise
affine background, is completely determined by the $A_\infty$
relations on $\FM_k(\C)$ (PVA Jacobi) and the $E_1$ relations on
$\Conf_k^{<}(\R)$.
\end{proof}

\begin{corollary}[The cubic identity; \ClaimStatusProvedHere]
\label{cor:cubic-identity}
At jet order $3$, the reflected weight identity reduces to:
\begin{equation}\label{eq:cubic-identity}
\delta_3^{\mathrm{refl}}(V)
\;=\;
-\,\delta_3^{\mathrm{alg}}(V).
\end{equation}
The left side is the self-image propagator integral of
Construction~\ref{cons:cubic-localization} (a lollipop graph with the
self-image kernel $K(z,t;\,z,-t)$ on the loop).  The right side is a
specific polynomial expression in the cubic structure constants
$c^{ij}_{k,\ell mn}$ and the affine structure constants
$f^{ij}_{k,\ell}$, computed by the unreflected one-loop Kontsevich
integral with a single cubic insertion.

The PVA Jacobi identity for $V$ determines
$\delta_3^{\mathrm{alg}}(V)$ uniquely.  The cubic reflected weight
identity (Conjecture~\ref{conj:cubic-rwi}) is equivalent to the
\emph{single analytic identity}~\eqref{eq:cubic-identity}.
\end{corollary}

\begin{proof}
Apply Theorem~\ref{thm:jacobi-reflection-decomposition} at $N=3$.
The algebraic part $\delta_3^{\mathrm{alg}}$ is computed by
Theorem~\ref{thm:algebraic-determination} as a sum over trees and
one-loop graphs with one cubic vertex in an affine background.  The
reflection part $\delta_3^{\mathrm{refl}}$ is the self-image integral
of Construction~\ref{cons:cubic-localization}.  The reflected weight
identity is the equality of these two expressions.
\end{proof}

\subsection{The doubling theorem}

The Jacobi-reflection decomposition identifies the remaining content
of the half-space programme: prove that the reflection integrals
match the algebraically determined expressions.  We now resolve this
by embedding the half-space problem into the doubled whole-space,
where the standard Kontsevich--FM machinery applies.

\begin{construction}[The doubled configuration space]
\label{cons:doubled-configuration}
Let $X=\C\times\R_{\ge 0}$ be the half-space and
$X_{\mathrm{dbl}}=\C\times\R$ the doubled space.  The reflection
$\sigma\colon X_{\mathrm{dbl}}\to X_{\mathrm{dbl}}$,
$\sigma(z,t)=(\bar z,-t)$, has $X$ as a fundamental domain and
fixed-point locus $\R\times\{0\}\subset X_{\mathrm{dbl}}$.

The \emph{method-of-images embedding}
\[
\iota\colon\Conf_n(X)
\hookrightarrow
\Conf_{2n}(X_{\mathrm{dbl}})^{\sigma}
\]
sends a configuration $(p_1,\ldots,p_n)$ with $p_i=(z_i,t_i)\in X$ to
the $\sigma$-invariant configuration
$(p_1,\ldots,p_n,\sigma(p_1),\ldots,\sigma(p_n))$.  Each original
point is paired with its image $\sigma(p_i)=(\bar z_i,-t_i)$.
\end{construction}

\begin{proposition}[Reflected propagator as pullback; \ClaimStatusProvedHere]
\label{prop:reflected-pullback}
Let $\omega_{ij}$ denote the propagator $1$-form on
$\Conf_2(X_{\mathrm{dbl}})$ for the HT theory on the
whole space.  Under the images embedding~$\iota$, the
reflected propagator form on $\Conf_2(X)$ is
\begin{equation}\label{eq:refl-as-pullback}
\omega_{ij}^{\mathrm{refl}}
\;=\;
\iota^*\!\bigl(\omega_{ij}+\omega_{i,\sigma(j)}\bigr),
\end{equation}
where $\omega_{i,\sigma(j)}$ is the propagator form between the
original point $p_i$ and the image point $\sigma(p_j)$.
\end{proposition}

\begin{proof}
By definition,
$K^{\mathrm{refl}}(z_i,t_i;\,z_j,t_j)
=K(z_i,t_i;\,z_j,t_j)+K(z_i,t_i;\,\bar z_j,-t_j)$.
The first term is $K(p_i,p_j)$; the second is $K(p_i,\sigma(p_j))$.
The corresponding $1$-forms on the configuration space are
$\omega_{ij}$ and $\omega_{i,\sigma(j)}$, and $\iota$ identifies
$\sigma(p_j)$ with the image point $(\bar z_j,-t_j)$ in the doubled
space.
\end{proof}

\begin{theorem}[The reflected weight identity from the doubling
principle; \ClaimStatusProvedHere]
\label{thm:doubling-rwi}
Assume hypotheses \textrm{(H1)--(H4)}.  Then the reflected weight
identity~\eqref{eq:reflected-weight-identity} holds for every
admissible graph $\Gamma$ on the half-space $X=\C\times\R_{\ge 0}$.
\end{theorem}

\begin{proof}
The proof has four steps: double, compactify, apply Stokes, restrict.

\medskip\noindent\textbf{Step~1: Doubled graph and configuration
space.}
Given an admissible graph $\Gamma$ on $n$ vertices in $X$, the
images embedding~$\iota$ produces a $\sigma$-invariant graph
$\widetilde\Gamma$ on $2n$ vertices in $X_{\mathrm{dbl}}$: for
each edge $e=(i,j)$ of $\Gamma$, the doubled graph
$\widetilde\Gamma$ has two edges $(i,j)$ and $(i,\sigma(j))$
(corresponding to the two terms in
$\omega^{\mathrm{refl}}=\omega+\omega_{\mathrm{img}}$).  The
$\sigma$-invariance of $\widetilde\Gamma$ follows from the
$\sigma$-invariance of the configuration.

\medskip\noindent\textbf{Step~2: FM compactification of the doubled
space.}
By hypothesis (H3), the propagator forms $\omega_{ij}$ extend to
logarithmic $1$-forms on the Fulton--MacPherson compactification
$\overline{\Conf}_{2n}(X_{\mathrm{dbl}})$.  The collision divisors
$D_S$ (for subsets $S\subseteq\{1,\ldots,2n\}$) are the standard
FM boundary strata.  The Arnold--Orlik--Solomon relations (H3) ensure
that codimension-$2$ corners cancel pairwise.

The FM compactification includes collision strata between original
points $(D_{\{i,j\}})$, between image points
$(D_{\{\sigma(i),\sigma(j)\}})$, and between an original point and
an image point $(D_{\{i,\sigma(j)\}})$).  The last type corresponds
precisely to the \emph{reflection strata} of
Construction~\ref{cons:half-space-colored-operad}: when $p_i$
collides with $\sigma(p_j)$ in $X_{\mathrm{dbl}}$, the half-space
projection sees $p_i$ approaching the image of $p_j$---a
reflection-type collision.

\medskip\noindent\textbf{Step~3: Stokes' theorem on the doubled
space.}
The wedge product
$\bigwedge_{e\in E(\widetilde\Gamma)}\omega_e$ is a closed
logarithmic form of top degree on the interior of
$\overline{\Conf}_{2n}(X_{\mathrm{dbl}})$ (each factor is closed by
(H2), and the FM compactification has normal-crossing boundary
divisors by (H3)).  Stokes' theorem gives
\begin{equation}\label{eq:stokes-doubled}
0
\;=\;
\int_{\overline{\Conf}_{2n}(X_{\mathrm{dbl}})}
d\Bigl(\bigwedge_{e\in E(\widetilde\Gamma)}\omega_e\Bigr)
\;=\;
\sum_{S}\int_{D_S}
\operatorname{Res}_{D_S}\Bigl(
\bigwedge_{e\in E(\widetilde\Gamma)}\omega_e\Bigr).
\end{equation}
This is the whole-space weight identity for the doubled graph
$\widetilde\Gamma$.

\medskip\noindent\textbf{Step~4: Restriction to the
$\sigma$-invariant locus.}
The $\sigma$-action on $\overline{\Conf}_{2n}(X_{\mathrm{dbl}})$
preserves the FM boundary structure: $\sigma$ maps $D_S$ to
$D_{\sigma(S)}$.  The $\sigma$-invariant
restriction of~\eqref{eq:stokes-doubled} is:
\begin{equation}\label{eq:stokes-restricted}
0
\;=\;
\sum_{S\;\sigma\text{-inv}}
\int_{D_S^{\sigma}}
\operatorname{Res}_{D_S}\Bigl(
\bigwedge_{e\in E(\widetilde\Gamma)}\omega_e\Bigr)\Bigg|_{\sigma}.
\end{equation}
The $\sigma$-invariant strata $D_S^{\sigma}$ decompose into
exactly the four types of
Construction~\ref{cons:half-space-colored-operad}:
\begin{itemize}
\item $S\subseteq\{1,\ldots,n\}$ (original points only): $D_S^{\sigma}$
  is a \textbf{bulk collision stratum} in $X$.
\item $S\subseteq\{1,\ldots,n\}$ restricted to the boundary
  $t_i=0$: $D_S^{\sigma}$ is a \textbf{boundary collision stratum}.
\item $S$ contains both original and image points from the same pair
  $\{i,\sigma(i)\}$: $D_S^{\sigma}$ is a \textbf{mixed
  bulk-boundary stratum}.
\item $S=\{i,\sigma(j)\}$ with $i\neq j$: $D_S^{\sigma}$ is a
  \textbf{reflection stratum} (original point $p_i$ collides with
  image point $\sigma(p_j)$).
\end{itemize}
By Proposition~\ref{prop:reflected-pullback}, the restricted form
$(\bigwedge_e\omega_e)|_{\sigma}=\bigwedge_e\omega_e^{\mathrm{refl}}$
is the reflected propagator form on $\Conf_n(X)$.  Therefore
the $\sigma$-restricted Stokes
identity~\eqref{eq:stokes-restricted} is exactly the reflected
weight identity~\eqref{eq:reflected-weight-identity}.
\end{proof}

\begin{theorem}[General half-space BV quantization; \ClaimStatusProvedHere]
\label{thm:general-half-space-bv}
Assume \textrm{(H1)--(H4)}.  Let $V$ be a freely generated Poisson
vertex algebra with inner holomorphic translation.  The Poisson--vertex
sigma model on $\C\times\R_{\ge 0}$ with Neumann boundary condition
admits a renormalized perturbative BV quantization.  The boundary
factorization algebra is a vertex algebra deforming~$V$.
\end{theorem}

\begin{proof}
By Theorem~\ref{thm:doubling-rwi}, the reflected weight identity
holds under (H1)--(H4).  By Theorem~\ref{thm:half-space-reduction}
(implication $(2)\Rightarrow(1)$), the boundary factorization algebra
is a vertex algebra.
\end{proof}

\begin{remark}[What the doubling proves]
The doubling argument shows that the half-space reflected weight
identity is not an independent analytic identity: it is the
$\sigma$-invariant restriction of the standard whole-space Stokes
identity.  The reflected propagator is not a new object: it is the
standard propagator on the doubled space, restricted to the
fundamental domain.  The reflection strata are not new strata: they
are standard collision strata in the doubled space, between an
original point and an image point.

The half-space geometry does not introduce new analytic content.  It
reorganizes existing content---the whole-space Kontsevich
machinery---through the symmetry $\sigma$.
\end{remark}

\begin{remark}[Relation to the Dolgushev theorem]
The doubling construction provides the geometric model for the
Dolgushev equivariant formality map.  The Dolgushev theorem
guarantees existence of a $\C^*$-equivariant deformation quantization
by abstract homotopy-transfer arguments.  The doubling construction
realizes this quantization concretely: the equivariant star product
is the restriction of the standard Kontsevich star product on
$X_{\mathrm{dbl}}$ to the $\sigma$-invariant sector.

The identification
\[
\text{Dolgushev equivariant formality}
\;=\;
\text{Kontsevich formality on }X_{\mathrm{dbl}}
\;\text{restricted to }\sigma\text{-invariants}
\]
makes the half-space BV quantization a consequence of the standard
Kontsevich formality theorem, applied to the doubled space and
restricted by symmetry.
\end{remark}

\section{Synthesis}

The half-space programme is complete under hypotheses (H1)--(H4):
\[
\text{PVA}
\;\xrightarrow{\text{PV sigma model}}
\text{half-space BV theory}
\;\xrightarrow{\text{doubling + Stokes}}
\text{reflected weight identity}
\;\xrightarrow{\text{recognition}}
\text{boundary vertex algebra}.
\]
The doubling construction (Theorem~\ref{thm:doubling-rwi}) resolves
the nonlinear frontier.  The jet-order scheme
(Theorem~\ref{thm:jet-inductive-rwi}), the Jacobi-reflection
decomposition (Theorem~\ref{thm:jacobi-reflection-decomposition}),
and the cubic identity (Corollary~\ref{cor:cubic-identity}) are now
consequences: each jet-order identity is the $\sigma$-restriction of
the whole-space Stokes identity at the corresponding filtration level.

The affine half-space theorem (Theorem~\ref{thm:affine-half-space-bv})
is the first case: for linear Poisson structures, the doubled
Kontsevich integral localizes on wheel graphs, giving one-loop
exactness and the level shift $K_{\mathrm{eff}}=k+\hbar\,\delta_\fg$.
For nonlinear Poisson structures, the doubled integral does not
localize on wheels, and higher-loop contributions are generically
nonzero---but the Stokes identity (and hence the reflected weight
identity) holds at all loop orders by the FM compactification
and the closedness of the propagator form.
